% ============================================================
% f2012_fujian_arts_q19_cuboid.tex
% 2012 福建卷（文科）第 19 题 长方体 ABCD-A1B1C1D1
% AB=AD=1，AA1=2，M 为棱 DD1 上一点
% 来源 PDF：普通高考/2012/2012福建文.pdf，印面 732 第 1 栏上部。
% ============================================================
\begin{tikzpicture}[>=Stealth, line width=0.7pt, font=\small, x=1cm, y=1cm]
  \coordinate (B) at (0.00,0.00);
  \coordinate (C) at (1.60,0.00);
  \coordinate (B1) at (0.00,2.55);
  \coordinate (C1) at (1.60,2.55);
  \coordinate (A) at (0.95,0.58);
  \coordinate (D) at (2.55,0.58);
  \coordinate (A1) at (0.95,3.13);
  \coordinate (D1) at (2.55,3.13);
  \coordinate (M) at (2.55,1.55);

  % 可见棱
  \draw (B) -- (C) -- (C1) -- (B1) -- cycle;
  \draw (B1) -- (A1) -- (D1) -- (C1);
  \draw (C) -- (D) -- (D1);

  % 不可见棱
  \draw[dashed] (B) -- (A) -- (D);
  \draw[dashed] (A) -- (A1);

  \fill (A1) circle (1.2pt) node[above left] {$A_1$};
  \fill (B1) circle (1.2pt) node[left] {$B_1$};
  \fill (C1) circle (1.2pt) node[right] {$C_1$};
  \fill (D1) circle (1.2pt) node[right] {$D_1$};
  \fill (A) circle (1.2pt) node[below left] {$A$};
  \fill (B) circle (1.2pt) node[below left] {$B$};
  \fill (C) circle (1.2pt) node[below right] {$C$};
  \fill (D) circle (1.2pt) node[right] {$D$};
  \fill (M) circle (1.2pt) node[right] {$M$};
\end{tikzpicture}
